\documentclass[9pt]{developercv} % Default font size, values from 8-12pt are recommended

% es-noquoting stops babel-spanish making " active, which would collide with the
% template's markup. Without babel this document was hyphenated with ENGLISH patterns.
\usepackage[spanish,es-noquoting]{babel}

% Metadata. pdflang=es-ES also gives parsers the signal to load a Spanish month table
% -- "Dic" is the one month abbreviation here that an English table does not resolve.
\hypersetup{
	pdftitle={Alin Trandafir -- Desarrollador Full-Stack Senior -- CV},
	pdfauthor={Alin Trandafir},
	pdfsubject={Desarrollador Full-Stack Senior -- TypeScript, React, Angular, React Native},
	pdfkeywords={TypeScript, JavaScript, React, React Native, Angular, Next.js, Node.js,
		Expo, EAS, Redux, RxJS, Firebase, Docker, Cypress, Jenkins, GitLab CI, CI/CD,
		Java, Symfony, Stripe, MapBox, i18next, internacionalizacion, i18n, l10n,
		Google Cloud, GCP, REST, SaaS, MongoDB, MySQL, SQL, NoSQL, Firestore,
		Keychain, Keystore, cifrado, CEO, Cofundador,
		autenticacion, MFA, Agile, Scrum,
		Desarrollador Full-Stack, Frontend, Movil, Fundador, Zaragoza, Espana, Remoto},
	pdflang={es-ES},
}

% "Nov 2024 -- Actualidad" needs 92.8pt; the 0.20 default gives 92.2pt and a \parbox
% overflows silently, running the date into the job title.
\setlength{\dateColWidth}{0.23\textwidth}

\begin{document}

\newlength{\namewidth}
\settowidth{\namewidth}{\HUGE\textbf{\MakeUppercase{Trandafir}}} % Size to the LONGER name

\begin{minipage}[t]{0.40\textwidth} % Left column: name and title
	\vspace{-\baselineskip} % Required for vertically aligning minipages

	\colorbox{cvink}{\parbox{\namewidth}{%
		\HUGE\textcolor{white}{\textbf{\MakeUppercase{Alin}}}\\[2pt]
		\HUGE\textcolor{white}{\textbf{\MakeUppercase{Trandafir}}}}}

	\vspace{6pt} % Vertical whitespace

	{\huge\raggedright Desarrollador Full-Stack\\Senior\par} % Career or current job title
\end{minipage}
\hfill % Automatic horizontal whitespace
\begin{minipage}[t]{0.56\textwidth} % Right column: contact details, stacked
	\vspace{-\baselineskip} % Required for vertically aligning minipages

	\icon{MapMarker}{12}{Zaragoza, España}\\
	\icon{Phone}{12}{Teléfono: +34 653 751 742}\\
	\icon{At}{12}{\href{mailto:alin\_trandafir@icloud.com}{alin\_trandafir@icloud.com}}\\
	\icon{Globe}{12}{Web: \href{https://end.works}{end.works}}\\
	\icon{Github}{12}{\href{https://github.com/ender-null}{github.com/ender-null}}\\
	\icon{Linkedin}{12}{\href{https://linkedin.com/in/alin-trandafir}{linkedin.com/in/alin-trandafir}}
\end{minipage}

\vspace{0.5cm} % Vertical whitespace

\cvsect{Perfil Profesional}

% "Perfil Profesional" rather than "¿Quién soy?": ATS segment a document by matching
% section headers against a dictionary, and "¿Quién soy?" matches nothing in any of them.

\textbf{Desarrollador full-stack senior} con \textasciitilde{}8 años de experiencia en \textbf{TypeScript}, \textbf{React}, \textbf{Angular} y \textbf{React Native}. He construido aplicaciones en producción y plataformas web escalables en grandes empresas, y desarrollo y opero las mías propias: \textbf{Zierzo} y \textbf{LoveLust} están publicadas en App Store y Google Play a través de \textbf{The Lovelust Company, SL}, la empresa que cofundé y dirijo como CEO, con LoveLust alcanzando los \textbf{16.000 usuarios activos mensuales}.
\\
Me siento cómodo \textbf{liderando la dirección técnica} y trabajando \textbf{en cualquier parte del stack} para obtener el mejor resultado posible. Me apasiona crear software que \textbf{facilite la vida de las personas} y resuelva problemas complejos -- eso es lo que me llevó a desarrollar mis propios productos.

\cvsect{Competencias Técnicas}


\textbf{Lenguajes:} TypeScript, JavaScript (ES6+), Java, PHP, HTML5, CSS3\\
\textbf{Frontend:} React, Angular, Next.js, Redux, RxJS, gestión de estado, librerías de componentes, diseño responsive\\
\textbf{Móvil:} React Native, Expo, EAS, Cordova, Keychain de iOS / Keystore de Android, publicación en App Store y Google Play\\
\textbf{Backend:} Node.js, APIs REST, Symfony\\
\textbf{Bases de datos:} MongoDB, MySQL, SQL, Firebase / Firestore\\
\textbf{Cloud y DevOps:} Google Cloud Platform (GCP), Docker, CI/CD, GitLab CI, Jenkins, Git\\
\textbf{Testing y Calidad:} Cypress, tests unitarios, tests end-to-end, cobertura de tests, code review\\
\textbf{Metodologías:} Agile, Scrum, internacionalización (i18n / l10n), autenticación y autorización, MFA, cifrado en reposo, SaaS, mentoría

\cvsect{Experiencia Laboral}

\begin{entrylist}
	% Ver nota en resume.tex: el tramo Nov 2024 - May 2026 se separa a propósito, para
	% que se lea como recorrido (construir -> tracción -> constituir la sociedad) y no
	% como un periodo suelto de freelance entre empleos.
	\entry
		{May 2026 -- Actualidad}
		{Cofundador y CEO}
		{The Lovelust Company, SL}
		{Cofundé la sociedad que opera \textbf{LoveLust} y \textbf{Zierzo}, y la dirijo como CEO sin dejar de ser su responsable técnico -- LoveLust llega a los \textbf{16.000 usuarios activos mensuales} con monetización por suscripción. Asumo ambos productos \textbf{de extremo a extremo y a escala mundial}: \textbf{arquitectura}, \textbf{base de datos}, analítica, compras integradas y publicación en ambas tiendas, además de \textbf{marketing}, finanzas y \textbf{cumplimiento normativo}. Represento a la empresa ante socios, prescriptores y organizaciones externas, y en eventos del sector. \\
		\texttt{React Native}\slashsep\texttt{Expo}\slashsep\texttt{TypeScript}\slashsep\texttt{RevenueCat}\slashsep\texttt{HealthKit/Health Connect}\slashsep\texttt{EAS}\slashsep\texttt{Firebase}\slashsep\texttt{i18next}}
	\entry
		{Nov 2024 -- May 2026}
		{Desarrollador de Software Independiente}
		{LoveLust \slashsep Zierzo}
		{Diseñé, desarrollé y publiqué dos productos desde cero hasta \textbf{App Store y Google Play} como único desarrollador: \textbf{cifrado de datos en local} con las claves custodiadas en el \textbf{Keychain de iOS} y el \textbf{Keystore de Android}, facturación por suscripción, integración con datos de salud, funcionalidades en tiempo real, analítica de uso y soporte multiidioma. Asumí todas las capas: \textbf{arquitectura}, \textbf{CI/CD}, \textbf{tests automatizados}, revisión en tiendas y despliegues progresivos. También desarrollé \textbf{sitios web para pequeños negocios locales}. La tracción de ambos productos llevó a cofundar \textbf{The Lovelust Company, SL} en mayo de 2026. \\
		\texttt{React Native}\slashsep\texttt{Expo}\slashsep\texttt{TypeScript}\slashsep\texttt{RevenueCat}\slashsep\texttt{EAS}\slashsep\texttt{Firebase}\slashsep\texttt{i18next}}
	\entry
		{Oct 2022 -- Nov 2024}
		{Ingeniero de Software}
		{Aimsun}
		{Lideré de extremo a extremo los módulos centrales de una \textbf{plataforma SaaS} -- autenticación, cuentas, suscripciones, licencias y MFA -- dentro de un equipo multifuncional. Desarrollé herramientas internas para \textbf{facturación}, \textbf{permisos} y administración de usuarios. Contribuí a bibliotecas de autenticación compartidas e interfaces de visualización de datos en tiempo real. \\
		\texttt{React}\slashsep\texttt{TypeScript}\slashsep\texttt{Next.js}\slashsep\texttt{i18next}\slashsep\texttt{MapBox}\slashsep\texttt{GitLab CI/CD}\slashsep\texttt{Cypress}}
	\entry
		{Dic 2021 -- Oct 2022}
		{Desarrollador de Software}
		{Viewnext}
		{Desarrollé una \textbf{plataforma de reservas de puntos de carga para vehículos eléctricos} para una importante empresa energética -- los usuarios podían consultar estaciones en un mapa interactivo y gestionar sus reservas. Contribuí a una \textbf{biblioteca de componentes compartida} con requisitos estrictos de calidad: cobertura de tests casi total y ausencia de code smells. Fui referente técnico de desarrolladores junior. \\
		\texttt{Angular}\slashsep\texttt{TypeScript}\slashsep\texttt{RxJS}\slashsep\texttt{Jenkins}\slashsep\texttt{GitLab CI/CD}}
	\entry
		{May 2021 -- Dic 2021}
		{Desarrollador de Software}
		{adidas}
		{Reemplacé herramientas internas legacy por un \textbf{editor de contenidos} moderno, reduciendo significativamente la fricción editorial. Introduje la \textbf{internacionalización} y optimicé los flujos de traducción en toda la plataforma. Trabajé en un entorno Agile junto a equipos de diseño, contenido e ingeniería. \\
		\texttt{React}\slashsep\texttt{TypeScript}\slashsep\texttt{Redux}\slashsep\texttt{i18n}}
	\entry
		{Dic 2019 -- May 2021}
		{Desarrollador de Software}
		{Telefónica / Eleven Paths}
		{Lideré la migración de una plataforma de gobernanza y gestión de riesgos a una arquitectura moderna con \textbf{Angular}. Entregué funcionalidades empresariales, mejoré la coherencia de la experiencia de usuario y mantuve servicios backend para informes y modelado de riesgos. \\
		\texttt{Angular}\slashsep\texttt{TypeScript}\slashsep\texttt{RxJS}\slashsep\texttt{Java}\slashsep\texttt{i18n}}
	\entry
		{Oct 2017 -- Dic 2019}
		{Desarrollador de Software}
		{Ingenia Sistemas Avanzados para el Transporte}
		{Desarrollé dos productos para empresas privadas de transporte de viajeros: una \textbf{plataforma de venta de billetes personalizable por operador} con descuentos, múltiples métodos de pago e integración mediante iframe; y una \textbf{herramienta de back-office interna} para gestionar estaciones, líneas, conductores y vehículos con generación de rutas mediante la \textbf{API de Google Maps}. También contribuí a una app móvil con llegadas en tiempo real. \\
		\texttt{Angular}\slashsep\texttt{TypeScript}\slashsep\texttt{RxJS}\slashsep\texttt{Cordova}\slashsep\texttt{Stripe}\slashsep\texttt{i18n}}
	\entry
		{Mar 2017 -- Sep 2017}
		{Desarrollador de Software}
		{Rananegra}
		{Desarrollé sitios web dinámicos con un CMS modular y componentes interactivos personalizados. Creé una herramienta de automatización de emails con facturación por suscripción. \\
		\texttt{Symfony}\slashsep\texttt{Angular}\slashsep\texttt{TypeScript}\slashsep\texttt{Stripe}}
\end{entrylist}

\cvsect{Certificaciones}

% Separadas de Formación y colocadas primero: son las credenciales más recientes y
% relevantes, y antes aparecían las últimas, por debajo de dos titulaciones de
% 2012-2016, en una lista cuyo orden era 2014, 2012, 2025, 2025.
%
% Emisores corregidos: "Formación San Miguel" y "Kings Corner English School" son las
% academias que prepararon los exámenes, NO los organismos certificadores.

\begin{entrylist}
	\shortentry
		{2025}
		{Google Cloud Professional Cloud Architect}
		{Google Cloud}
		{}

	\shortentry
		{2025}
		{Oxford Test of English -- Advanced (C1)}
		{Oxford University Press}
		{}
\end{entrylist}

\cvsect{Formación}

% CORRECCIÓN IMPORTANTE: DAM es un Ciclo Formativo de Grado Superior -> "Técnico
% Superior"; SMR es un Ciclo Formativo de Grado Medio -> "Técnico". En España "Grado"
% designa un título universitario, así que la redacción anterior ("Grado en...")
% atribuía una titulación que no se posee -- algo que cualquier recruiter español
% detecta de inmediato. No hay ninguna desventaja en el título correcto: un CFGS en DAM
% es una vía plenamente reconocida, y con 8 años de experiencia nadie evalúa esto.

\begin{entrylist}
	\shortentry
		{2014 -- 2016}
		{Técnico Superior en Desarrollo de Aplicaciones Multiplataforma}
		{CPIFP Los Enlaces}
		{}

	\shortentry
		{2012 -- 2014}
		{Técnico en Sistemas Microinformáticos y Redes}
		{CPIFP Los Enlaces}
		{}
\end{entrylist}

% Pie a DOS columnas, no tres. La columna "Soft Skills" desaparece: su encabezado
% seguía en INGLÉS en el CV español (deriva heredada de copiar el fichero), y sus
% contenidos eran o bien no verificables o bien universales. Lo sustantivo (Agile/Scrum,
% mentoría, liderazgo) vive ahora en Competencias Técnicas > Metodologías.

\begin{minipage}[t]{0.42\textwidth} % Left column
	\vspace{-\baselineskip} % Required for vertically aligning minipages

	\cvsect{Idiomas}
	\textbf{Español} -- nativo\\
	\textbf{Rumano} -- nativo (bilingüe de origen)\\
	\textbf{Inglés} -- C1 (Oxford Test of English, 2025)
\end{minipage}
\hfill % Automatic horizontal whitespace
\begin{minipage}[t]{0.52\textwidth} % Right column
	\vspace{-\baselineskip} % Required for vertically aligning minipages

	\cvsect{Voluntariado}
	\textbf{Fundador y Presidente}, Asociación Juvenil Brovrashgka -- asociación juvenil cultural sin ánimo de lucro registrada, dedicada a la cultura japonesa y la comunidad creativa. Gestioné el \textbf{presupuesto} anual y las solicitudes de subvención, coordiné a los \textbf{voluntarios} y organicé eventos y talleres públicos de forma recurrente.
\end{minipage}

\end{document}
